% from ega0-20-fr.tex
\oldpage[0\textsubscript{IV-1}]{136}
\begin{remark}[20.5.15]
\label{0.20.5.15-fr}
Soit $u:A\to B$ un homomorphisme \emph{surjectif} d'anneaux; alors, pour
tout homomorphisme d'anneaux $v:B\to C$, l'homomorphisme canonique
\begin{equation}
\tag{20.5.15.1}
\label{0.20.5.15.1-fr}
  u_{C/B/A}:\Omega^1_{C/A}\longrightarrow\Omega^1_{C/B}
\end{equation}
est \emph{bijectif}; cela résulte en effet de la suite exacte (20.5.7.1),
puisque $\Omega^1_{B/A}=0$ (20.4.12).
\end{remark}

\subsection{Modules d'imperfection et homomorphismes caractéristiques}
\label{subsection:0.20.6-fr}
\label{0.20.6-fr}

\begin{definition}[20.6.1]
\label{0.20.6.1-fr}
Étant donnés deux homomorphismes d'anneaux $u:A\to B$, $v:B\to C$, on
appelle \emph{module d'imperfection de la $B$-algèbre $C$ relativement à
$A$}, et on note $\Upsilon_{C/B/A}$, le $C$-module noyau de
l'homomorphisme
$v_{C/B/A}:\Omega^1_{B/A}\otimes_B C\to\Omega^1_{C/A}$.
On a donc par définition (cf. (20.5.7)) la suite exacte
\begin{equation}
\tag{20.6.1.1}
\label{0.20.6.1.1-fr}
  0\longrightarrow\Upsilon_{C/B/A}\longrightarrow
  \Omega^1_{B/A}\otimes_B C\xrightarrow{v_{C/B/A}}
  \Omega^1_{C/A}\xrightarrow{u_{C/B/A}}\Omega^1_{C/B}
  \longrightarrow0.
\end{equation}
\end{definition}

Lorsque $A=\mathbb{Z}$ (les modules $\Omega^1_{B/A}$ et $\Omega^1_{C/A}$
étant donc les modules de différentielles « absolues » $\Omega^1_B$ et
$\Omega^1_C$), on écrit $\Upsilon_{C/B}$ au lieu de
$\Upsilon_{C/B/\mathbb{Z}}$; lorsque $B$ et $C$ sont des algèbres sur un
corps premier $P$, on a $\Upsilon_{C/B/P}=\Upsilon_{C/B}$.

Soient $R$, $S$ des parties multiplicatives de $B$ et $C$ respectivement,
telles que l'image de $R$ soit contenue dans $S$. Il résulte alors de la
suite exacte (20.6.1.1) et de (20.5.9) que l'on a
\begin{equation}
\tag{20.6.1.2}
\label{0.20.6.1.2-fr}
  \Upsilon_{S^{-1}C/R^{-1}B/A}=S^{-1}\Upsilon_{C/B/A}.
\end{equation}

\begin{proposition}[20.6.2]
\label{0.20.6.2-fr}
Si $C$ est une $B$-algèbre formellement lisse relativement à $A$ (et en
particulier si $C$ est une $B$-algèbre formellement lisse), on a
$\Upsilon_{C/B/A}=0$.
\end{proposition}

\begin{proof}
Cela résulte de (20.5.7, (ii)).
\end{proof}

\begin{proposition}[20.6.3]
\label{0.20.6.3-fr}
Soient $k$ un corps, $K$ une extension de $k$. Pour que $K$ soit une
extension séparable de $k$, il faut et il suffit que $\Upsilon_{K/k}=0$
(autrement dit, l'homomorphisme canonique
$\Omega^1_k\otimes_k K\to\Omega^1_K$ est injectif ou encore (20.4.8),
toute dérivation de $k$ dans $K$ se prolonge en une dérivation de $K$
dans lui-même).
\end{proposition}

\begin{proof}
En effet, il est équivalent de dire que $K$ est séparable sur $k$ ou une
$k$-algèbre formellement lisse (19.6.1); d'autre part, si $P$ est le
corps premier de $k$, $K$ est séparable sur $P$, donc une $P$-algèbre
formellement lisse, et il revient donc au même de dire que $K$ est une
$k$-algèbre formellement lisse ou une $k$-algèbre formellement lisse
relativement à $P$. Enfin, dire que $K$ est une $k$-algèbre formellement
lisse relativement à $P$ équivaut d'après (20.5.7, (ii)) à dire que
l'homomorphisme
$v_{K/k/P}:\Omega^1_k\otimes_k K\to\Omega^1_K$ est inversible à gauche;
mais comme $K$ est un corps, cette dernière condition équivaut à dire
que le noyau de $v_{K/k/P}$, c'est-à-dire $\Upsilon_{K/k}$, est nul.
\end{proof}

\begin{env}[20.6.4]
\label{0.20.6.4-fr}
Considérons un diagramme commutatif
\begin{equation}
\tag{20.6.4.1}
\label{0.20.6.4.1-fr}
\xymatrix{
  A' \ar[r]^{u'} & B' \ar[r]^{v'} & C' \\
  A \ar[u]^f \ar[r]_u & B \ar[u]^g \ar[r]_v & C \ar[u]^h
}
\end{equation}

